\documentclass[11pt]{ctexart}
\usepackage[a4paper,margin=1in]{geometry}
\usepackage{longtable,booktabs}
\title{Geant4-Agent Casebank Regression Report (EN)}
\author{Automated Pipeline}
\date{2026-03-08\_113601}
\begin{document}
\maketitle
\section{Unified Pass/Fail Table}
\begin{longtable}{p{0.06\linewidth}p{0.08\linewidth}p{0.07\linewidth}p{0.08\linewidth}p{0.07\linewidth}p{0.07\linewidth}p{0.07\linewidth}p{0.07\linewidth}p{0.09\linewidth}p{0.06\linewidth}p{0.21\linewidth}}
\toprule
ID & Domain & Style & Completeness & ExpInit & ActInit & ExpFinal & ActFinal & Missing(I->F) & Result & Fail Reason \\
\midrule
MTM01 & medical & fuzzy & missing & FAIL & FAIL & PASS & PASS & 12->0 & PASS & - \\
MTM02 & aerospace & precise & missing & FAIL & FAIL & PASS & PASS & 2->0 & PASS & - \\
MTM03 & industrial\_ndt & precise & missing & FAIL & FAIL & PASS & PASS & 12->0 & PASS & - \\
MTM04 & physics & precise & missing & FAIL & FAIL & PASS & PASS & 8->0 & PASS & - \\
MTM05 & nuclear\_engineering & fuzzy & missing & FAIL & FAIL & PASS & PASS & 9->0 & PASS & - \\
MTM06 & security\_screening & fuzzy & missing & FAIL & FAIL & PASS & PASS & 11->0 & PASS & - \\
MTM07 & semiconductor & precise & missing & FAIL & FAIL & PASS & PASS & 11->0 & PASS & - \\
MTM08 & education & fuzzy & missing & FAIL & FAIL & PASS & PASS & 11->0 & PASS & - \\
MTM09 & environmental\_radiation & fuzzy & missing & FAIL & FAIL & PASS & PASS & 3->0 & PASS & - \\
MTM10 & high\_energy\_physics & precise & missing & FAIL & FAIL & PASS & PASS & 10->0 & PASS & - \\
MTM11 & materials\_update & precise & full & PASS & PASS & PASS & PASS & 0->0 & PASS & - \\
MTM12 & overwrite\_reject & precise & full & PASS & PASS & PASS & PASS & 0->0 & PASS & - \\
MTM13 & medical\_detector & fuzzy & missing & FAIL & FAIL & PASS & PASS & 13->0 & PASS & - \\
MTM14 & industrial\_panel & precise & missing & FAIL & FAIL & PASS & PASS & 15->0 & FAIL & expected\_final\_mismatch:materials.selected\_materials.0 \\
MTM15 & shielding\_stack & precise & missing & FAIL & FAIL & PASS & PASS & 14->0 & PASS & - \\
MTM16 & nested\_target & precise & missing & FAIL & FAIL & PASS & PASS & 9->0 & FAIL & expected\_final\_mismatch:source.type \\
MTM17 & coaxial\_shield & fuzzy & missing & FAIL & FAIL & PASS & PASS & 10->0 & PASS & - \\
MTM18 & boolean\_cavity & precise & missing & FAIL & FAIL & PASS & FAIL & 17->1 & FAIL & final\_complete\_mismatch, expected\_final\_mismatch:materials.selected\_materials.0 \\
MTM19 & ring\_complete & precise & full & PASS & PASS & PASS & PASS & 0->0 & PASS & - \\
MTM20 & grid\_complete & precise & full & PASS & FAIL & PASS & FAIL & 1->1 & FAIL & initial\_complete\_mismatch, final\_complete\_mismatch, expected\_final\_mismatch:geometry.structure, expected\_final\_mismatch:materials.selected\_materials.0 \\
MTM21 & layered\_detector & fuzzy & missing & FAIL & FAIL & PASS & PASS & 14->0 & PASS & - \\
MTM22 & embedded\_sensor & precise & full & PASS & PASS & PASS & PASS & 0->0 & PASS & - \\
MTM23 & layered\_shield\_complete & precise & full & PASS & PASS & PASS & PASS & 0->0 & PASS & - \\
MTM24 & boolean\_complete & precise & full & PASS & FAIL & PASS & FAIL & 1->1 & FAIL & initial\_complete\_mismatch, final\_complete\_mismatch, expected\_final\_mismatch:materials.selected\_materials.0, expected\_final\_mismatch:source.type \\
\bottomrule
\end{longtable}
\section{Baseline Comparison}
\begin{itemize}
\item pass\_rate=0.792
\item initial\_alignment\_rate=0.917
\item final\_alignment\_rate=0.875
\item avg\_turns\_per\_case=2.458
\item missing\_reduction\_rate=1.000
\item confirm\_apply\_success\_rate=1.000
\item rollback\_after\_confirm\_rate=0.000
\item internal\_leak\_turn\_rate=0.051
\item repeated\_assistant\_case\_rate=0.000
\item stale\_confirm\_case\_rate=0.000
\item delta\_pass\_rate=-0.158
\item regressions=none
\item improvements=none
\end{itemize}
\appendix
\section{Appendix: Full Dialogues and Naturalized Internal Chains}
\subsection{Case MTM01 (medical)}
\textbf{Dialogue:}
\begin{itemize}
\item User: 我先要一个用于医学成像的初步辐照场景，细节还没定，你先告诉我还缺什么。
\item Assistant: 好的，我们来看看目前的配置进度：

**已完成：**
*   已经确定了使用“盒体”作为几何结构。

**待补充：**
1.  **几何**：需要您提供盒体在 X、Y、Z 三个方向的具体尺寸。
2.  **材料**：需要指定这个盒体体积使用什么材料。
3.  **源**：需要选择放射源的类型（例如伽马源）、粒子种类和能量。

请随时告诉我这些信息，我们继续完善配置。
\item User: 采用 300 mm x 300 mm x 300 mm 的水箱，材料 ?；源用束流 gamma，能量 6 MeV，位置 (0,0,-400) mm，方向 +z；物理列表 QBBC；输出 root。
\item Assistant: 我检测到这会覆盖已经确认的值。待修改内容：几何结构类型：unknown -> single\_box；盒体 x 方向尺寸：None -> 300.0。如果要应用这些修改，请直接回复“确认”；如果要保留原值，请直接说“保持原值”。
\item User: 确认
\item Assistant: 好的，配置已完成！目前我们已经确认了盒体的几何结构类型以及它在X、Y、Z三个方向上的尺寸。
\end{itemize}
\textbf{Internal Chain (Natural Language):}
\begin{itemize}
\item Turn 1: action=summarize\_progress, phase=geometry, complete=False.
\item Missing fields: geometry.params.module\_x, geometry.params.module\_y, geometry.params.module\_z, materials.volume\_material\_map, source.type, source.particle, source.energy, source.position, source.direction, physics.physics\_list, output.format, output.path.
\item LLM used=True; fallback=none; backend=llm\_semantic\_frame+runtime\_semantic.
\item Turn 2: action=confirm\_overwrite, phase=geometry, complete=False.
\item Missing fields: geometry.params.module\_x, geometry.params.module\_y, geometry.params.module\_z.
\item LLM used=True; fallback=none; backend=llm\_slot\_frame+runtime\_semantic.
\item Turn 3: action=finalize, phase=materials, complete=True.
\item Missing fields: none.
\item LLM used=True; fallback=none; backend=llm\_slot\_frame+runtime\_semantic+confirmed\_pending\_overwrite.
\end{itemize}
\subsection{Case MTM02 (aerospace)}
\textbf{Dialogue:}
\begin{itemize}
\item User: I need a first-pass proton irradiation setup for a silicon device; keep it simple for now.
\item Assistant: Geometry: updated geometry structure, box size along x, box size along y. Materials: updated selected materials, volume-to-material binding. Source: updated source type, particle type, source energy. Physics: updated physics list. Output: still needs output format, output path.
\item User: Use a silicon box 50 mm x 50 mm x 5 mm, beam source proton 100 MeV from (0,0,-100) mm along +z, physics FTFP\_BERT, output json.
\item Assistant: This would change a value that's already been set.

Pending change:
- Box size along x: 1000.0 → 50.0
- Box size along y: 1000.0 → 50.0

Reply "confirm" to apply this change, or "keep existing" to keep the current values.
\item User: confirm
\item Assistant: Your configuration is complete. The following items have been confirmed: source type, and the box size along the x, y, and z axes.
\end{itemize}
\textbf{Internal Chain (Natural Language):}
\begin{itemize}
\item Turn 1: action=summarize\_progress, phase=materials, complete=False.
\item Missing fields: output.format, output.path.
\item LLM used=True; fallback=none; backend=llm\_slot\_frame+runtime\_semantic.
\item Turn 2: action=confirm\_overwrite, phase=source, complete=False.
\item Missing fields: none.
\item LLM used=True; fallback=none; backend=llm\_slot\_frame+runtime\_semantic.
\item Turn 3: action=finalize, phase=physics, complete=True.
\item Missing fields: none.
\item LLM used=True; fallback=none; backend=llm\_semantic\_frame+runtime\_semantic+confirmed\_pending\_overwrite.
\end{itemize}
\subsection{Case MTM03 (industrial\_ndt)}
\textbf{Dialogue:}
\begin{itemize}
\item User: 我要做一个 gamma 穿透铁靶的小测试，但我还没把完整参数说完。
\item Assistant: 目前我们已经确定了使用铁作为材料，并选择伽马粒子进行模拟。

接下来还需要你补充几何结构的具体类型和尺寸（例如盒体的长、宽、高），以及指定体积与材料的绑定关系。
\item User: 请设为半径 40 mm、半长 80 mm 的铁圆柱，材料 ?；点源 gamma，能量 2 MeV，位置 (0,0,-200) mm，方向 +z；物理列表 Shielding；输出 csv。
\item Assistant: 配置已完成。当前已确认的参数包括：体积与材料绑定、输出路径、输出格式和物理列表。
\end{itemize}
\textbf{Internal Chain (Natural Language):}
\begin{itemize}
\item Turn 1: action=summarize\_progress, phase=geometry, complete=False.
\item Missing fields: geometry.structure, geometry.params.module\_x, geometry.params.module\_y, geometry.params.module\_z, materials.volume\_material\_map, source.type, source.energy, source.position, source.direction, physics.physics\_list, output.format, output.path.
\item LLM used=True; fallback=none; backend=llm\_slot\_frame+runtime\_semantic.
\item Turn 2: action=finalize, phase=materials, complete=True.
\item Missing fields: none.
\item LLM used=True; fallback=none; backend=llm\_slot\_frame+runtime\_semantic.
\end{itemize}
\subsection{Case MTM04 (physics)}
\textbf{Dialogue:}
\begin{itemize}
\item User: Gamma point source 1 MeV at (0,0,-100) toward +z; please build the rest.
\item Assistant: Geometry: still needs geometry structure, box size along x. Materials: still needs volume-to-material binding. Source: updated source type, particle type, source energy. Physics: still needs physics list. Output: still needs output format, output path.
\item User: Set copper box 1 m x 1 m x 1 m, physics FTFP\_BERT, output hdf5.
\item Assistant: Your configuration is complete. The following items are confirmed: volume-to-material binding, output path, output format, and physics list.
\end{itemize}
\textbf{Internal Chain (Natural Language):}
\begin{itemize}
\item Turn 1: action=summarize\_progress, phase=geometry, complete=False.
\item Missing fields: geometry.structure, geometry.params.module\_x, geometry.params.module\_y, geometry.params.module\_z, materials.volume\_material\_map, physics.physics\_list, output.format, output.path.
\item LLM used=True; fallback=none; backend=llm\_slot\_frame+runtime\_semantic.
\item Turn 2: action=finalize, phase=materials, complete=True.
\item Missing fields: none.
\item LLM used=True; fallback=none; backend=llm\_slot\_frame+runtime\_semantic.
\end{itemize}
\subsection{Case MTM05 (nuclear\_engineering)}
\textbf{Dialogue:}
\begin{itemize}
\item User: 我想做一个中子屏蔽的概念验证，现在先给一个缺参开头。
\item Assistant: 目前进度：几何部分已选定结构类型，材料、粒子类型和物理列表也已同步。接下来还需要补充盒体的具体尺寸、源的类型与能量，以及输出设置。
\item User: 使用 200 mm x 200 mm x 200 mm 的铅块，材料 ?；点源 neutron，能量 14 MeV，位置 (0,0,-50) mm，方向 +z；物理列表 QGSP\_BERT\_HP；输出 xml。
\item Assistant: 我注意到这次设置会覆盖一些已经确认过的参数。具体来说：

- 几何结构类型：从“未知”改为“单个盒体”
- 已选材料：从“铜”改为“铅”

请确认你的选择：
• 如果要应用这些修改，请直接回复“确认”
• 如果要保留原来的值，请直接说“保持原值”
\item User: 确认
\item Assistant: 配置已完成。当前已确认的参数包括：体积与材料绑定、盒体在 x、y、z 三个方向的尺寸。
\end{itemize}
\textbf{Internal Chain (Natural Language):}
\begin{itemize}
\item Turn 1: action=summarize\_progress, phase=geometry, complete=False.
\item Missing fields: geometry.params.module\_x, geometry.params.module\_y, geometry.params.module\_z, source.type, source.energy, source.position, source.direction, output.format, output.path.
\item LLM used=True; fallback=none; backend=llm\_slot\_frame+runtime\_semantic.
\item Turn 2: action=confirm\_overwrite, phase=geometry, complete=False.
\item Missing fields: geometry.params.module\_x, geometry.params.module\_y, geometry.params.module\_z.
\item LLM used=True; fallback=none; backend=llm\_slot\_frame+runtime\_semantic.
\item Turn 3: action=finalize, phase=materials, complete=True.
\item Missing fields: none.
\item LLM used=True; fallback=none; backend=llm\_semantic\_frame+runtime\_semantic+confirmed\_pending\_overwrite.
\end{itemize}
\subsection{Case MTM06 (security\_screening)}
\textbf{Dialogue:}
\begin{itemize}
\item User: Need a quick cargo-screening toy setup with forward photons, but I have not fixed the target shape, material, source energy, or output format yet.
\item Assistant: Here's what we have so far:

**Geometry:** You've chosen the structure. Next, please provide the box dimensions (size along X and Y).
**Materials:** We still need to assign a material to the volume.
**Source:** You've selected photons. Next, we need the source type and its energy.
**Physics:** A physics list is still required.
**Output:** An output format and file path are still needed.
\item User: Use steel cylinder target radius 60 mm half-length 120 mm, beam gamma 3 MeV from (0,0,-500) mm to +z, physics Shielding, output root.
\item Assistant: This would overwrite an already confirmed value. Pending change: geometry structure: unknown -> single\_tubs; box size along x: None -> 1000.0. Reply 'confirm' to apply it, or say 'keep existing' to keep the current value.
\item User: confirm
\item Assistant: Configuration complete. The following items are now confirmed: the volume's material binding, the cylinder's half-length and radius, and the box size along the Z-axis.
\end{itemize}
\textbf{Internal Chain (Natural Language):}
\begin{itemize}
\item Turn 1: action=summarize\_progress, phase=geometry, complete=False.
\item Missing fields: geometry.params.module\_x, geometry.params.module\_y, geometry.params.module\_z, materials.volume\_material\_map, source.type, source.energy, source.position, source.direction, physics.physics\_list, output.format, output.path.
\item LLM used=True; fallback=none; backend=llm\_slot\_frame+runtime\_semantic.
\item Turn 2: action=confirm\_overwrite, phase=geometry, complete=False.
\item Missing fields: geometry.params.module\_x, geometry.params.module\_y, geometry.params.module\_z.
\item LLM used=True; fallback=none; backend=llm\_slot\_frame+runtime\_semantic.
\item Turn 3: action=finalize, phase=materials, complete=True.
\item Missing fields: none.
\item LLM used=True; fallback=none; backend=llm\_semantic\_frame+runtime\_semantic+confirmed\_pending\_overwrite.
\end{itemize}
\subsection{Case MTM07 (semiconductor)}
\textbf{Dialogue:}
\begin{itemize}
\item User: 我想搭一个硅薄片的中子照射测试，但先只给一个简述。
\item Assistant: 目前我们已经确定了使用中子源和硅材料。接下来还需要你提供几何结构的具体类型和尺寸，以及将材料绑定到体积上。请告诉我你希望构建哪种几何形状（比如盒子、圆柱等）以及它的尺寸参数。
\item User: 使用 20 mm x 20 mm x 1 mm 的硅片，材料 ?；各向同性 neutron，能量 2 MeV，位置 (0,0,0) mm；物理列表 QGSP\_BERT\_HP；输出 json。
\item Assistant: 我注意到这将会覆盖一个已经确认的设置：源类型将从“beam”改为“isotropic”。

请确认你的选择：
- 如果要应用这个修改，请直接回复“确认”。
- 如果要保留原来的“beam”设置，请直接回复“保持原值”。
\item User: 确认
\item Assistant: 配置已完成。当前已确认的设置包括：源类型、输出路径、输出格式和物理列表。
\end{itemize}
\textbf{Internal Chain (Natural Language):}
\begin{itemize}
\item Turn 1: action=summarize\_progress, phase=geometry, complete=False.
\item Missing fields: geometry.structure, geometry.params.module\_x, geometry.params.module\_y, geometry.params.module\_z, materials.volume\_material\_map, source.energy, source.position, source.direction, physics.physics\_list, output.format, output.path.
\item LLM used=True; fallback=none; backend=llm\_semantic\_frame+runtime\_semantic.
\item Turn 2: action=confirm\_overwrite, phase=materials, complete=False.
\item Missing fields: none.
\item LLM used=True; fallback=none; backend=llm\_slot\_frame+runtime\_semantic.
\item Turn 3: action=finalize, phase=source, complete=True.
\item Missing fields: none.
\item LLM used=True; fallback=none; backend=llm\_semantic\_frame+runtime\_semantic+confirmed\_pending\_overwrite.
\end{itemize}
\subsection{Case MTM08 (education)}
\textbf{Dialogue:}
\begin{itemize}
\item User: For classroom demo, I want gamma attenuation in copper; you can ask only essential missing parts.
\item Assistant: Geometry: still needs geometry structure, box size along x. Materials: updated selected materials; still needs volume-to-material binding. Source: updated particle type; still needs source type, source energy. Physics: updated physics list. Output: still needs output format, output path.
\item User: Okay, copper box 100 mm x 100 mm x 100 mm, point gamma 1 MeV at (0,0,-200) mm, direction +z, physics FTFP\_BERT, output csv.
\item Assistant: You've already confirmed the physics list as "Shielding." Changing it to "FTFP\_BERT" would overwrite that.

Please reply **'confirm'** to apply the change to FTFP\_BERT, or **'keep existing'** to keep the current Shielding list.
\item User: confirm
\item Assistant: Your configuration is complete. The physics list, output path, output format, and source direction have been confirmed.
\end{itemize}
\textbf{Internal Chain (Natural Language):}
\begin{itemize}
\item Turn 1: action=summarize\_progress, phase=geometry, complete=False.
\item Missing fields: geometry.structure, geometry.params.module\_x, geometry.params.module\_y, geometry.params.module\_z, materials.volume\_material\_map, source.type, source.energy, source.position, source.direction, output.format, output.path.
\item LLM used=True; fallback=none; backend=llm\_slot\_frame+runtime\_semantic.
\item Turn 2: action=confirm\_overwrite, phase=materials, complete=False.
\item Missing fields: none.
\item LLM used=True; fallback=none; backend=llm\_slot\_frame+runtime\_semantic.
\item Turn 3: action=finalize, phase=source, complete=True.
\item Missing fields: none.
\item LLM used=True; fallback=none; backend=llm\_semantic\_frame+runtime\_semantic+confirmed\_pending\_overwrite.
\end{itemize}
\subsection{Case MTM09 (environmental\_radiation)}
\textbf{Dialogue:}
\begin{itemize}
\item User: 我想先搭一个空气体积中的低能 gamma 传播场景，细节稍后补。
\item Assistant: 目前已完成几何、材料和源的配置。接下来请补充物理列表、输出格式和输出路径。
\item User: 设为 500 mm x 500 mm x 500 mm 的空气盒，材料 ??；点源 gamma，能量 0.5 MeV，位置 (0,0,-300) mm，方向 +z；物理列表 QBBC；输出 root。
\item Assistant: 我注意到这将会覆盖之前已经确认的参数。具体修改如下：

- 盒体 x 方向尺寸：从 1000.0 改为 500.0
- 盒体 y 方向尺寸：从 1000.0 改为 500.0

如果你确认应用这些修改，请直接回复“确认”；如果想保留原来的数值，请直接说“保持原值”。
\item User: 确认
\item Assistant: 配置已完成。当前已确认的参数包括：源位置、源能量、源类型以及盒体 z 方向尺寸。
\end{itemize}
\textbf{Internal Chain (Natural Language):}
\begin{itemize}
\item Turn 1: action=summarize\_progress, phase=materials, complete=False.
\item Missing fields: physics.physics\_list, output.format, output.path.
\item LLM used=True; fallback=none; backend=llm\_slot\_frame+runtime\_semantic.
\item Turn 2: action=confirm\_overwrite, phase=source, complete=False.
\item Missing fields: none.
\item LLM used=True; fallback=none; backend=llm\_slot\_frame+runtime\_semantic.
\item Turn 3: action=finalize, phase=physics, complete=True.
\item Missing fields: none.
\item LLM used=True; fallback=none; backend=llm\_semantic\_frame+runtime\_semantic+confirmed\_pending\_overwrite.
\end{itemize}
\subsection{Case MTM10 (high\_energy\_physics)}
\textbf{Dialogue:}
\begin{itemize}
\item User: Set geometry as a sphere target; I have not decided radius or source details yet.
\item Assistant: You've set the geometry to a sphere target. We still need the sphere's radius.

We also need to assign a material to that volume, and define your particle source (its type, particle, and energy).

Finally, we'll need to select a physics list and specify your output preferences.
\item User: Sphere target radius 80 mm, lead material, electron point source 5 MeV at (0,0,-200) mm toward +z, physics FTFP\_BERT, output hdf5.
\item Assistant: Your configuration is complete. The following items have been confirmed: the material assignment for the volume, the output path and format, and the physics list.
\end{itemize}
\textbf{Internal Chain (Natural Language):}
\begin{itemize}
\item Turn 1: action=summarize\_progress, phase=geometry, complete=False.
\item Missing fields: geometry.params.child\_rmax, materials.volume\_material\_map, source.type, source.particle, source.energy, source.position, source.direction, physics.physics\_list, output.format, output.path.
\item LLM used=True; fallback=none; backend=llm\_slot\_frame+runtime\_semantic.
\item Turn 2: action=finalize, phase=materials, complete=True.
\item Missing fields: none.
\item LLM used=True; fallback=none; backend=llm\_slot\_frame+runtime\_semantic.
\end{itemize}
\subsection{Case MTM11 (materials\_update)}
\textbf{Dialogue:}
\begin{itemize}
\item User: Create a complete setup: 100 mm x 100 mm x 100 mm copper box, point gamma 1 MeV at (0,0,-200) mm toward +z, physics QBBC, output json.
\item Assistant: Your configuration is complete. The volume-to-material binding, output path, and output format have been confirmed.
\item User: Change the target material to aluminum and change the output format to root.
\item Assistant: You're about to overwrite two confirmed settings:
- Selected materials: from copper to aluminum.
- Output format: from JSON to ROOT.

Please reply **'confirm'** to apply these changes, or **'keep existing'** to keep the current values.
\item User: confirm
\item Assistant: Your configuration is complete. The volume-to-material binding, output path, output format, and selected materials have been confirmed.
\end{itemize}
\textbf{Internal Chain (Natural Language):}
\begin{itemize}
\item Turn 1: action=finalize, phase=materials, complete=True.
\item Missing fields: none.
\item LLM used=True; fallback=none; backend=llm\_slot\_frame+runtime\_semantic.
\item Turn 2: action=confirm\_overwrite, phase=source, complete=False.
\item Missing fields: none.
\item LLM used=True; fallback=none; backend=llm\_slot\_frame+runtime\_semantic.
\item Turn 3: action=finalize, phase=physics, complete=True.
\item Missing fields: none.
\item LLM used=True; fallback=none; backend=llm\_semantic\_frame+runtime\_semantic+confirmed\_pending\_overwrite.
\end{itemize}
\subsection{Case MTM12 (overwrite\_reject)}
\textbf{Dialogue:}
\begin{itemize}
\item User: 请创建一个完整配置：100 mm x 100 mm x 100 mm 的铜盒，材料 ?；点源 gamma，能量 1 MeV，位置 (0,0,-200) mm，方向 +z；物理列表 QBBC；输出 json。
\item Assistant: 配置已完成。当前已确认的设置包括：体积与材料绑定、输出路径和输出格式。
\item User: 把材料改成 ?，并把输出格式改成 root。
\item Assistant: 我注意到这次修改会覆盖已确认的设置。具体来说：输出格式将从“json”改为“root”。

请确认你的选择：
- 如果要应用这个修改，请直接回复“确认”。
- 如果要保留原来的“json”格式，请直接回复“保持原值”。
\item User: 保持原值
\item Assistant: 好的，输出格式没有更改，仍然使用原来的“json”格式。
\end{itemize}
\textbf{Internal Chain (Natural Language):}
\begin{itemize}
\item Turn 1: action=finalize, phase=materials, complete=True.
\item Missing fields: none.
\item LLM used=True; fallback=none; backend=llm\_slot\_frame+runtime\_semantic.
\item Turn 2: action=confirm\_overwrite, phase=source, complete=False.
\item Missing fields: geometry.ask.boolean.solid\_a\_size, geometry.ask.boolean.solid\_b\_size.
\item LLM used=True; fallback=none; backend=llm\_slot\_frame+runtime\_semantic.
\item Turn 3: action=reject\_overwrite, phase=physics, complete=True.
\item Missing fields: none.
\item LLM used=True; fallback=none; backend=llm\_semantic\_frame+runtime\_semantic+rejected\_pending\_overwrite.
\end{itemize}
\subsection{Case MTM13 (medical\_detector)}
\textbf{Dialogue:}
\begin{itemize}
\item User: I want a ring of detector modules around a small source region, but I have not fixed the ring radius or module size yet.
\item Assistant: I've updated your geometry to a ring structure. To continue, I need:

*   **Geometry:** The size and number of detector modules for the ring.
*   **Materials:** Which material to assign to each volume.
*   **Source:** The type of particle, its energy, and how the source is defined.
*   **Physics \& Output:** Your choices for the physics list and how to save the results.
\item User: Use a ring of 12 modules, each 8 mm x 10 mm x 2 mm, radius 40 mm, clearance 1 mm, material water; point gamma source 0.662 MeV at (0,0,0) mm; physics QBBC; output json.
\item Assistant: This would overwrite an already confirmed geometry value.

**Pending change:**
*   **Module size:** From (7.83, 4.72, 6.86) mm to (8, 10, 2) mm.
*   **Ring:** From 4 modules at 67.49 mm radius (0.19 mm clearance) to 12 modules at 40 mm radius (1 mm clearance).

Reply **'confirm'** to apply this change, or say **'keep existing'** to keep the current value.
\item User: confirm
\item Assistant: Your configuration is complete. The following geometry items have been confirmed: ring clearance, ring radius, ring module count, and ring module size.
\end{itemize}
\textbf{Internal Chain (Natural Language):}
\begin{itemize}
\item Turn 1: action=summarize\_progress, phase=materials, complete=False.
\item Missing fields: materials.volume\_material\_map, source.type, source.particle, source.energy, source.position, source.direction, physics.physics\_list, output.format, output.path, geometry.ask.ring.module\_size, geometry.ask.ring.count, geometry.ask.ring.radius, geometry.ask.ring.clearance.
\item LLM used=True; fallback=none; backend=llm\_slot\_frame+runtime\_semantic.
\item Turn 2: action=confirm\_overwrite, phase=source, complete=False.
\item Missing fields: none.
\item LLM used=True; fallback=none; backend=llm\_slot\_frame+runtime\_semantic.
\item Turn 3: action=finalize, phase=physics, complete=True.
\item Missing fields: none.
\item LLM used=True; fallback=none; backend=llm\_semantic\_frame+runtime\_semantic+confirmed\_pending\_overwrite.
\end{itemize}
\subsection{Case MTM14 (industrial\_panel)}
\textbf{Dialogue:}
\begin{itemize}
\item User: 我想做一个二维阵列探测板，但 pitch 和材料还没定。
\item Assistant: 好的，我们已确认了您要构建的二维阵列探测板几何结构。

**当前进度：**
*   ✅ 几何部分：已确定使用“盒体”作为基本结构。
*   ⏳ 材料、源、物理和输出部分的信息尚待补充。

**接下来，请您提供以下信息以继续配置：**
1.  每个探测单元的尺寸（阵列模块尺寸）。
2.  探测单元在X和Y方向上的数量。
3.  为探测板指定的材料（体积与材料绑定）。
4.  模拟使用的放射源类型、粒子种类及其能量。
\item User: ?? 4 x 3 ?????????? 10 mm x 10 mm x 2 mm?pitch\_x 12 mm?pitch\_y 14 mm??? 1 mm??? ???? gamma 0.511 MeV??? (0,0,-80) mm??? +z????? QBBC??? root?
\item Assistant: 我检测到这会覆盖已经确认的值。待修改内容：几何结构类型：grid -> single\_box；盒体 x 方向尺寸：None -> 10.0。如果要应用这些修改，请直接回复“确认”；如果要保留原值，请直接说“保持原值”。
\item User: 确认
\item Assistant: 好的，配置已完成！

当前我们已经确认了以下参数：
*   y 方向 pitch
*   x 方向 pitch
*   盒体 z 方向尺寸
*   盒体 y 方向尺寸
*   盒体 x 方向尺寸
\end{itemize}
\textbf{Internal Chain (Natural Language):}
\begin{itemize}
\item Turn 1: action=summarize\_progress, phase=materials, complete=False.
\item Missing fields: materials.volume\_material\_map, source.type, source.particle, source.energy, source.position, source.direction, physics.physics\_list, output.format, output.path, geometry.ask.grid.module\_size, geometry.ask.grid.count\_x, geometry.ask.grid.count\_y, geometry.ask.grid.pitch\_x, geometry.ask.grid.pitch\_y, geometry.ask.grid.clearance.
\item LLM used=True; fallback=none; backend=llm\_semantic\_frame+runtime\_semantic.
\item Turn 2: action=confirm\_overwrite, phase=source, complete=False.
\item Missing fields: none.
\item LLM used=True; fallback=none; backend=llm\_slot\_frame+runtime\_semantic.
\item Turn 3: action=finalize, phase=physics, complete=True.
\item Missing fields: none.
\item LLM used=True; fallback=none; backend=llm\_slot\_frame+runtime\_semantic+confirmed\_pending\_overwrite.
\end{itemize}
\subsection{Case MTM15 (shielding\_stack)}
\textbf{Dialogue:}
\begin{itemize}
\item User: Build a stacked absorber sandwich along z; I will provide the layer sizes next.
\item Assistant: You've set up a stacked absorber sandwich geometry along the z-axis. Here's what's complete and what's still needed:

**Completed:**
*   **Geometry:** The basic structure is set.

**Still Needed:**
*   **Geometry:** The footprint and thicknesses for each stack layer.
*   **Materials:** Which material to assign to each volume.
*   **Source:** The source type, particle type, and its energy.
*   **Physics:** The physics list to use.
*   **Output:** The format and file path for your results.
\item User: Use stacked layers with footprint 20 mm x 20 mm, thicknesses 2 mm, 4 mm, 6 mm, layer clearance 0.5 mm, outer box 30 mm x 30 mm x 20 mm, container clearance 1 mm, material aluminum; beam gamma 1.25 MeV from (0,0,-120) mm along +z; physics Shielding; output csv.
\item Assistant: You are about to overwrite a geometry configuration that was already confirmed.

**Pending Change:**
*   **Stack Footprint:** 16.33 mm x 10.28 mm → 20.0 mm x 20.0 mm
*   **Layer Thicknesses:** [6.68, 1.19, 5.59] mm → [2.0, 4.0, 6.0] mm
*   **Layer Clearance:** 0.73 mm → 0.5 mm
*   **Container Size:** 18.07 mm x 45.52 mm x 16.66 mm → 30.0 mm x 30.0 mm x 20.0 mm
*   **Container Clearance:** 0.87 mm → 1.0 mm

Reply **'confirm'** to apply this change, or **'keep existing'** to keep the current values.
\item User: confirm
\item Assistant: Configuration complete. The following geometry items are confirmed: stack container clearance, stack container size, stack layer clearance, and stack layer thicknesses.
\end{itemize}
\textbf{Internal Chain (Natural Language):}
\begin{itemize}
\item Turn 1: action=summarize\_progress, phase=materials, complete=False.
\item Missing fields: materials.volume\_material\_map, source.type, source.particle, source.energy, source.position, source.direction, physics.physics\_list, output.format, output.path, geometry.ask.stack.footprint, geometry.ask.stack.thicknesses, geometry.ask.stack.layer\_clearance, geometry.ask.stack.parent\_size, geometry.ask.stack.nest\_clearance.
\item LLM used=True; fallback=none; backend=llm\_slot\_frame+runtime\_semantic.
\item Turn 2: action=confirm\_overwrite, phase=source, complete=False.
\item Missing fields: none.
\item LLM used=True; fallback=none; backend=llm\_slot\_frame+runtime\_semantic.
\item Turn 3: action=finalize, phase=physics, complete=True.
\item Missing fields: none.
\item LLM used=True; fallback=none; backend=llm\_semantic\_frame+runtime\_semantic+confirmed\_pending\_overwrite.
\end{itemize}
\subsection{Case MTM16 (nested\_target)}
\textbf{Dialogue:}
\begin{itemize}
\item User: 我想在一个盒子里嵌一个圆柱靶，先别急着补全。
\item Assistant: 好的，我们已同步了以下信息：
- **几何**：结构类型（盒中嵌圆柱），以及盒体的尺寸。
- **源**：能量、位置和方向。

接下来，还需要您补充以下内容才能完成配置：
- **材料**：为盒体和圆柱体分别指定材料。
- **源**：选择粒子类型（如质子、电子等）。
- **物理**：选择物理列表。
- **输出**：确定输出格式和保存路径。
\item User: 外盒 120 mm x 120 mm x 120 mm，内嵌铜圆柱半径 20 mm、半长 40 mm，嵌套间隙 2 mm；材料 ?；束流 gamma 1 MeV，位置 (0,0,-150) mm，方向 +z；物理列表 QBBC；输出 json。
\item Assistant: 我注意到这次更新会覆盖一些已经确认过的参数。具体来说：

1.  **源位置**：从 (0, 0, -100) mm 改为 (0, 0, -150) mm。
2.  **几何结构**：外层盒体和内嵌圆柱的尺寸、间隙都更新为您刚指定的数值。

请问您希望如何处理？
- 如果确认应用这些修改，请直接回复 **“确认”**。
- 如果希望保留原来的参数值，请直接回复 **“保持原值”**。
\item User: 确认
\item Assistant: 配置已完成。当前已确认的参数包括：源位置、嵌套间隙、内嵌圆柱半长和内嵌圆柱半径。
\end{itemize}
\textbf{Internal Chain (Natural Language):}
\begin{itemize}
\item Turn 1: action=summarize\_progress, phase=materials, complete=False.
\item Missing fields: materials.volume\_material\_map, source.particle, physics.physics\_list, output.format, output.path, geometry.ask.nest.parent\_size, geometry.ask.nest.child\_radius, geometry.ask.nest.child\_half\_length, geometry.ask.nest.clearance.
\item LLM used=True; fallback=none; backend=llm\_slot\_frame+runtime\_semantic.
\item Turn 2: action=confirm\_overwrite, phase=source, complete=False.
\item Missing fields: none.
\item LLM used=True; fallback=none; backend=llm\_slot\_frame+runtime\_semantic.
\item Turn 3: action=finalize, phase=physics, complete=True.
\item Missing fields: none.
\item LLM used=True; fallback=none; backend=llm\_semantic\_frame+runtime\_semantic+confirmed\_pending\_overwrite.
\end{itemize}
\subsection{Case MTM17 (coaxial\_shield)}
\textbf{Dialogue:}
\begin{itemize}
\item User: I need a shell-style shielding toy model, but the shell thicknesses are still open.
\item Assistant: Geometry: updated geometry structure; still needs shell inner radius, shell thicknesses. Materials: still needs volume-to-material binding. Source: updated source energy, source position, source direction; still needs particle type. Physics: updated physics list. Output: still needs output format, output path.
\item User: Use a shell with inner radius 15 mm, thicknesses 5 mm and 8 mm, half length 40 mm, nested core radius 10 mm, nested core half length 30 mm, clearance 1 mm, material lead; isotropic neutron source 2 MeV at (0,0,0) mm; physics Shielding; output hdf5.
\item Assistant: This would overwrite an already confirmed value. The pending changes are:
- Source type: from 'point' to 'isotropic'
- Source energy: from 1.0 to 2.0

Please reply **'confirm'** to apply these changes, or **'keep existing'** to keep the current values.
\item User: confirm
\item Assistant: Configuration complete. The following items are confirmed: source position, source energy, source type, and shell clearance.
\end{itemize}
\textbf{Internal Chain (Natural Language):}
\begin{itemize}
\item Turn 1: action=summarize\_progress, phase=materials, complete=False.
\item Missing fields: materials.volume\_material\_map, source.particle, output.format, output.path, geometry.ask.shell.inner\_radius, geometry.ask.shell.thicknesses, geometry.ask.shell.half\_length, geometry.ask.shell.child\_radius, geometry.ask.shell.child\_half\_length, geometry.ask.shell.clearance.
\item LLM used=True; fallback=none; backend=llm\_slot\_frame+runtime\_semantic.
\item Turn 2: action=confirm\_overwrite, phase=source, complete=False.
\item Missing fields: none.
\item LLM used=True; fallback=none; backend=llm\_slot\_frame+runtime\_semantic.
\item Turn 3: action=finalize, phase=physics, complete=True.
\item Missing fields: none.
\item LLM used=True; fallback=none; backend=llm\_semantic\_frame+runtime\_semantic+confirmed\_pending\_overwrite.
\end{itemize}
\subsection{Case MTM18 (boolean\_cavity)}
\textbf{Dialogue:}
\begin{itemize}
\item User: 我想做一个挖空的盒体几何，但先只给你意图。
\item Assistant: 好的，我们已确认了您要创建“挖空盒体”的几何结构，并记录了盒体的尺寸。

接下来，为了完成这个布尔减运算的几何，还需要您提供被挖去的那个实体（Solid A）的尺寸，包括它的X、Y和Z方向长度。
\item User: ???????A ? 40 mm x 40 mm x 40 mm ???B ? 20 mm x 20 mm x 20 mm ???? A ??? B??? ???? gamma 0.3 MeV??? (0,0,-60) mm??? +z????? QBBC??? xml?
\item Assistant: 我注意到这将会覆盖已经确认的设置。具体来说：

- 几何骨架会从“布尔减运算”改为“布尔并运算”。
- 几何程序中的尺寸和操作也会相应更新。

如果您确认应用这些修改，请直接回复“确认”；如果您希望保持原来的设置不变，请直接说“保持原值”。
\item User: 确认
\item Assistant: 好的，我们目前已经确认了以下设置：

*   **几何**：两个实体的尺寸。
*   **源**：方向、位置和能量。
*   **物理**：物理列表。
*   **输出**：文件路径和格式。

接下来，还需要您为几何体中的各个**体积指定材料**，这样我们就能完成全部配置了。
\end{itemize}
\textbf{Internal Chain (Natural Language):}
\begin{itemize}
\item Turn 1: action=summarize\_progress, phase=geometry, complete=False.
\item Missing fields: geometry.params.bool\_a\_x, geometry.params.bool\_a\_y, geometry.params.bool\_a\_z, geometry.params.bool\_b\_x, geometry.params.bool\_b\_y, geometry.params.bool\_b\_z, materials.volume\_material\_map, source.type, source.particle, source.energy, source.position, source.direction, physics.physics\_list, output.format, output.path, geometry.ask.boolean.solid\_a\_size, geometry.ask.boolean.solid\_b\_size.
\item LLM used=True; fallback=none; backend=llm\_slot\_frame+runtime\_semantic.
\item Turn 2: action=confirm\_overwrite, phase=materials, complete=False.
\item Missing fields: materials.volume\_material\_map.
\item LLM used=True; fallback=none; backend=llm\_slot\_frame+runtime\_semantic.
\item Turn 3: action=summarize\_progress, phase=materials, complete=False.
\item Missing fields: materials.volume\_material\_map.
\item LLM used=True; fallback=none; backend=llm\_semantic\_frame+runtime\_semantic+confirmed\_pending\_overwrite.
\end{itemize}
\subsection{Case MTM19 (ring\_complete)}
\textbf{Dialogue:}
\begin{itemize}
\item User: Create a complete ring setup: 16 modules, each 6 mm x 6 mm x 2 mm, radius 32 mm, clearance 0.5 mm, material silicon; point gamma source 0.511 MeV at (0,0,0) mm; physics QBBC; output json.
\item Assistant: Your configuration is complete. I have set up the ring geometry with 16 silicon modules, a point gamma source at 0.511 MeV, and the QBBC physics list for JSON output.

The following items are confirmed: volume-to-material binding, output path, clearance, and radius.
\end{itemize}
\textbf{Internal Chain (Natural Language):}
\begin{itemize}
\item Turn 1: action=finalize, phase=materials, complete=True.
\item Missing fields: none.
\item LLM used=True; fallback=none; backend=llm\_slot\_frame+runtime\_semantic.
\end{itemize}
\subsection{Case MTM20 (grid\_complete)}
\textbf{Dialogue:}
\begin{itemize}
\item User: ???????????3 x 3 ????????? 12 mm x 12 mm x 3 mm?pitch\_x 15 mm?pitch\_y 15 mm??? 1 mm??? ?????? gamma 0.662 MeV??? (0,0,-90) mm??? +z????? QBBC??? root?
\item Assistant: 目前我们已经确认了几何、源、物理和输出的设置，材料部分还需要补充体积与材料绑定的信息。
\end{itemize}
\textbf{Internal Chain (Natural Language):}
\begin{itemize}
\item Turn 1: action=summarize\_progress, phase=materials, complete=False.
\item Missing fields: materials.volume\_material\_map.
\item LLM used=True; fallback=none; backend=llm\_slot\_frame+runtime\_semantic.
\end{itemize}
\subsection{Case MTM21 (layered\_detector)}
\textbf{Dialogue:}
\begin{itemize}
\item User: I want a stacked detector slab, but the individual layer thicknesses are not fixed yet.
\item Assistant: You've set up a stacked detector geometry. Here's what we have and what's still needed:

**Completed:**
*   **Geometry:** The basic structure (a stacked slab) is configured. The overall box size is confirmed.

**Still Needed:**
*   **Geometry:** The footprint (area) and thicknesses for each stack layer.
*   **Materials:** Which material to assign to each volume.
*   **Source:** The source type, particle type, and its energy.
*   **Physics:** The physics list to use.
*   **Output:** The format and file path for saving results.
\item User: Use stack footprint 25 mm x 25 mm, thicknesses 1 mm, 2 mm, 1 mm, layer clearance 0.2 mm, outer box 30 mm x 30 mm x 12 mm, container clearance 0.5 mm, material silicon; beam electron 5 MeV from (0,0,-50) mm along +z; physics FTFP\_BERT; output json.
\item Assistant: This would overwrite an already confirmed geometry configuration.

**Pending Change:**
*   **Stack footprint:** 16.33 mm x 10.28 mm → 25.0 mm x 25.0 mm
*   **Stack layer thicknesses:** [6.68, 1.19, 5.59] mm → [1.0, 2.0, 1.0] mm
*   **Stack layer clearance:** 0.73 mm → 0.2 mm
*   **Outer box size:** 18.07 mm x 45.52 mm x 16.66 mm → 30.0 mm x 30.0 mm x 12.0 mm
*   **Container clearance:** 0.87 mm → 0.5 mm

Reply **'confirm'** to apply these new values, or **'keep existing'** to keep the current geometry.
\item User: confirm
\item Assistant: Your configuration is complete. The following geometry items are confirmed: stack container clearance, stack container size, stack layer clearance, and stack layer thicknesses.
\end{itemize}
\textbf{Internal Chain (Natural Language):}
\begin{itemize}
\item Turn 1: action=summarize\_progress, phase=materials, complete=False.
\item Missing fields: materials.volume\_material\_map, source.type, source.particle, source.energy, source.position, source.direction, physics.physics\_list, output.format, output.path, geometry.ask.stack.footprint, geometry.ask.stack.thicknesses, geometry.ask.stack.layer\_clearance, geometry.ask.stack.parent\_size, geometry.ask.stack.nest\_clearance.
\item LLM used=True; fallback=none; backend=llm\_semantic\_frame+runtime\_semantic.
\item Turn 2: action=confirm\_overwrite, phase=source, complete=False.
\item Missing fields: none.
\item LLM used=True; fallback=none; backend=llm\_slot\_frame+runtime\_semantic.
\item Turn 3: action=finalize, phase=physics, complete=True.
\item Missing fields: none.
\item LLM used=True; fallback=none; backend=llm\_semantic\_frame+runtime\_semantic+confirmed\_pending\_overwrite.
\end{itemize}
\subsection{Case MTM22 (embedded\_sensor)}
\textbf{Dialogue:}
\begin{itemize}
\item User: 请创建完整的嵌套配置：外盒 80 mm x 80 mm x 80 mm，内嵌铅圆柱半径 15 mm、半长 25 mm，嵌套间隙 1 mm；材料 ?；点源 gamma 1 MeV，位置 (0,0,-100) mm，方向 +z；物理列表 Shielding；输出 csv。
\item Assistant: 配置已完成！当前已确认的设置包括：体积与材料绑定、输出路径、嵌套间隙（clearance）以及圆柱半径。
\end{itemize}
\textbf{Internal Chain (Natural Language):}
\begin{itemize}
\item Turn 1: action=finalize, phase=materials, complete=True.
\item Missing fields: none.
\item LLM used=True; fallback=none; backend=llm\_slot\_frame+runtime\_semantic.
\end{itemize}
\subsection{Case MTM23 (layered\_shield\_complete)}
\textbf{Dialogue:}
\begin{itemize}
\item User: Create a complete shell configuration: inner radius 12 mm, thicknesses 4 mm and 6 mm, half length 35 mm, nested core radius 8 mm, nested core half length 25 mm, clearance 1 mm, material iron; isotropic gamma source 0.8 MeV at (0,0,0) mm; physics QBBC; output hdf5.
\item Assistant: Your configuration is complete. The following items have been confirmed: source direction, and the box dimensions along the X, Y, and Z axes.
\end{itemize}
\textbf{Internal Chain (Natural Language):}
\begin{itemize}
\item Turn 1: action=finalize, phase=materials, complete=True.
\item Missing fields: none.
\item LLM used=True; fallback=none; backend=llm\_slot\_frame+runtime\_semantic.
\end{itemize}
\subsection{Case MTM24 (boolean\_complete)}
\textbf{Dialogue:}
\begin{itemize}
\item User: ?????????????????A ? 30 mm x 30 mm x 20 mm ???B ? 10 mm x 10 mm x 40 mm ???? A ? union??? ???? gamma 0.5 MeV??? (0,0,-70) mm??? +z????? QBBC??? json?
\item Assistant: 好的，目前配置进度如下：

**已完成设置：**
- **几何**：已确认结构类型与部分尺寸。
- **源**：已设置好类型、粒子、能量、位置和方向。
- **物理**：已选定物理列表。
- **输出**：已确定格式和路径。

**仍需补充：**
- **材料**：需要为几何体指定材料（体积与材料绑定）。

请提供材料信息，我们即可完成全部配置。
\end{itemize}
\textbf{Internal Chain (Natural Language):}
\begin{itemize}
\item Turn 1: action=summarize\_progress, phase=materials, complete=False.
\item Missing fields: materials.volume\_material\_map.
\item LLM used=True; fallback=none; backend=llm\_slot\_frame+runtime\_semantic.
\end{itemize}
\end{document}